\scanpage{22}

 gerbe est alors une \inserted{\unsure{Gr-}}catégorie $A'$-linéaire, canoniquement déterminée.

\deleted{Supposons que \deleted{tt $X\in$} \inserted{pour} tout $X\in\operatorname{Ob}C$, et considérons $O=\operatorname{End}(X)$, et
\[
N_{\mathrm{red}}:\unsure{k}\ O^*\longrightarrow A^*
\]}

Nous désirons, sous certaines conditions, en définir un objet inversible canonique, en établissant un système transitif d'isomorphismes entre les $\det(X_{A'})$ $(X\in\operatorname{Ob}C)$. Pour ceci il suffit de définir un isom canonique
\[
(\alpha)\qquad \det P\simeq A
\]
pour tt. module à droite inv. $P$ sur un $O=\operatorname{End}(X)$ $(X\in\operatorname{Ob}C)$, fonctoriel en $P$ (ce qui suppose que
\[
N_{\mathrm{red}}:O'^*\longrightarrow A^*\quad\text{est à valeurs dans }1\}
\]
$O'$ étant le commutant de $O$ dans $P$ ....) et multiplicatif si $\unsure{{}_{O''}Q_{O'}P_O}\,\unsure{\otimes}$
\[
N_{\mathrm{red}}(Q\otimes_{O'}P,O)\simeq N_{\mathrm{red}}(Q,O')\otimes N_{\mathrm{red}}(P,O)
\]
doit être compatible avec les isom \unsure{ass.}. Cela marche si on suppose
\margin{[posant $\Phi=N_{\mathrm{red}}(E^*)\subset K^*$, où $O=\operatorname{End}(X)$, $E=O\otimes_A K$]}
\[
\left\{\begin{array}{l}
\operatorname{Pic}(A)=0\\
\deleted{N_{\mathrm{red}}(E^*)\cap N_{\mathrm{red}}(O)}\\
\Phi\cap A^*=\{1\},\quad \Phi.A^*=K^*
\end{array}\right.
\]
Ces conditions sont vérifiées \deleted{\unsure{pour si} $A=\deleted{\unsure{\mathbb Q}}\inserted{\mathbb Z}$, \illegible{1}} \deleted{\unsure{corps de nb. tot. réel}}, si $A=\mathbb Z$, $E_K(=E_{\mathbb R})$ \unsure{pas décomposé} sur $\mathbb R$. Donc dans ce cas, on a un objet inv. canonique de degré $n$ de $\operatorname{Rep}(G_{A'})$, de sorte que le lien \unsure{structural} de la gerbe $G_{A'}$

\mspage{21}
